\documentclass[10pt,a4paper]{amsart}
\usepackage[margin=19mm]{geometry}
\usepackage{amsmath,amssymb,mathtools}
\usepackage[scr=rsfs,cal=euler]{mathalfa}
\usepackage{xfrac}
\usepackage[unicode,hidelinks,bookmarks=false]{hyperref}
\newcommand{\ie}{i.e.\ }
\newcommand{\cf}{cf.\ }
\newcommand{\resp}{respectively\ }
\newcommand{\Subsection}{\subsection}
\providecommand{\nobreakdash}{\nobreak\hbox{-}}
\setlength{\emergencystretch}{2em}
\multlinegap0pt
\usepackage{bm}
\usepackage{bbm}
\usepackage{dsfont}
\usepackage{xspace}
\usepackage{cancel}
\usepackage{mathtools}
\usepackage{amsthm}
\newtheorem{theorem}{Theorem}
\newtheorem{definition}[theorem]{Definition}
\newtheorem{lemma}[theorem]{Lemma}
\title[Geodesic X-ray transform and streaking artifacts on simple s]{Geodesic X-ray transform and streaking artifacts on simple surfaces or on spaces of constant curvature}
\author{Hiroyuki Chihara}
\date{Source: arXiv:2402.06899v4 (CC BY 4.0)}
\begin{document}
\maketitle

\noindent\emph{Source and licence.} Geodesic X-ray transform and streaking artifacts on simple surfaces or on spaces of constant curvature. Version arXiv:2402.06899v4 (CC BY 4.0), distributed under the Creative Commons Attribution 4.0 International licence, as recorded in the arXiv licence metadata for this exact version and shown on the abstract page https://arxiv.org/abs/2402.06899v4. Full rights notice: licenses/arxiv-2402.06899-CC-BY-4.0.txt.\par\smallskip

\noindent\textit{Editorial scope.} Counted unit: the Lemma whose source label is \texttt{theorem:4-1} in the article (source0.tex lines 2221--2239, source label \texttt{theorem:4-1}), delivered together with the article's own complete original proof of it (source0.tex lines 2242--2497, one \texttt{proof} environment of 256 lines). No mathematics is written, repaired or re-derived: statement and proof are the article's own text line for line. Required context retained uncounted: theorem:geodesicxraytransform (lines 570--591); equation:etav (lines 1067--1080). Every definition, display and label of those lines is retained unchanged. Omitted from this package and not redistributed: 1--569, 592--1066, 1081--2220, 2240--2241, 2498--7655. Nothing inside a retained excerpt is omitted or altered: every retained body is delivered exactly as the article's lines are written, so no omission receipt of any kind accompanies this package. Citations: the retained excerpts carry no citation command at all, so no bibliography entry of the article is transcribed and no reference is redistributed. Reference adaptation: none was needed and none was made; every reference inside the delivered unit resolves inside it, so no unresolved reference or citation is produced. Printed slips kept literal: no printed word, symbol, hypothesis or number of the counted statement or of its proof was altered. Layout and class: the article's own class is \texttt{amsart} with packages including amsmath, amstext, amsbsy, amsopn, upref, amsthm, amsfonts, amssymb, mathrsfs, mathtools, times, cite, bm, graphicx; the wrapper is the lane's pinned \texttt{amsart} editorial wrapper at 10pt on A4, which restates the theorem-like environments the retained text needs and adds the article's own macro definitions that the retained text needs (none), with extra wrapper packages (bm, bbm, dsfont, xspace, cancel, mathtools). The delivered numbering is the unit's own. Assets: no retained span contains a figure environment, a graphics-inclusion command, a PDF-page inclusion, a table or a TikZ picture; no image file is copied into this package, the package declares no asset dependency and no image file is redistributed with it. Licence: version arXiv:2402.06899v4 of this article is distributed under the Creative Commons Attribution 4.0 International licence, as recorded in the arXiv licence metadata for this exact version and shown on the abstract page https://arxiv.org/abs/2402.06899v4. A full rights notice is delivered at licenses/arxiv-2402.06899-CC-BY-4.0.txt. Characters: every retained excerpt is pure ASCII, so the delivered engine, XeLaTeX with its default Latin Modern text font, reports no missing character.
\begin{definition}[{\bf Geodesic X-ray transform of a compactly supported half density}] 
\label{theorem:geodesicxraytransform}
The geodesic X-ray transform on $(M,g)$ is defined for $f \in C^\infty_0(\Omega_{M^\text{int}}^{1/2})$ as a half density on $\partial_-S(M)$ by 
%
%
\begin{equation}
\mathcal{X}[f](w)
:=
\left(
\int_0^{\tau_+(w)}
\frac{f}{\lvert{dv_g}\rvert^{1/2}}
\bigl(\dot{\gamma}_w(s)\bigr)
ds
\right)
\lvert{d\mu(w)}\rvert^{1/2},
\quad
w \in \partial_-S(M). 
\label{equation:geodesicxraytransform} 
\end{equation}
%
%
\end{definition} 

\begin{equation}
\eta(v)
=
\eta\bigl(D\pi\vert_v\ddot{\gamma}_v(0)\bigr)
=
D\pi\vert^T_v\eta\bigl(\ddot{\gamma}_v(0)\bigr)
=
DF\vert^T_v\xi\bigl(\ddot{\gamma}_v(0)\bigr)
=
\xi\bigl(DF\vert_v\ddot{\gamma}_v(0)\bigr)
=
0. 
\label{equation:etav}
\end{equation}

\begin{lemma}
\label{theorem:4-1}
Let $j=1,\dotsc,J$. 
Then the composition $C_\mathcal{X}{\circ}N^\ast(\partial{D_j})$  is transversal. 
Moreover if we set 
%
%
$$
\Sigma_j
:=
\pi_{\partial_-S(M)}
\bigl(C_\mathcal{X}{\circ}N^\ast(\partial{D_j})\bigr), 
\quad
j=1,\dotsc,J,
$$
%
%
then $\operatorname{codim}(\Sigma_j)=1$ and $N^\ast(\Sigma_j)=C_\mathcal{X}{\circ}N^\ast(\partial{D_j})$. 
\end{lemma} 

\begin{proof}
We first show that $C_\mathcal{X}{\circ}N^\ast(D_j)$  is transversal. 
Set 
%
%
\begin{align*}
  \mathcal{Z}
& =
  T^\ast\bigl(\partial_-S(M)\bigr)
  \times
  T^\ast(M^\text{int})
  \times
  T^\ast(M^\text{int}),
\\
  \mathcal{D}
& =
  T^\ast\bigl(\partial_-S(M)\bigr)
  \times
  \Delta\bigl(T^\ast(M^\text{int})\bigr)
\end{align*}
%
%
for short. We have 
$\operatorname{codim}(\mathcal{D})=\operatorname{dim}\bigl(T^\ast(M^\text{int})\bigr)=2n$. 
In view of Theorem~\ref{theorem:geodesicxraytransform} and \eqref{equation:etav}, we have  
%
%
\begin{align*}
  C_\mathcal{X}{\times}N^\ast({\partial}D_j)
& =
  \bigl\{
  (\xi,\eta,\zeta) 
  \in 
  T^\ast\bigl(\partial_-S(M)\bigr){\times}T^\ast(M^\text{int}){\times}N^\ast(\partial{D_j}): 
\\
& \qquad
  \exists v \in S(M^\text{int})\ \text{s.t.}\  
  \eta \in T^\ast_{\pi(v)}(M^\text{int})\setminus\{0\}, 
  \eta(v)=0, 
\\
& \qquad 
  \xi \in T^\ast_{F(v)}\bigl(\partial_-S(M)\bigr)\setminus\{0\} 
  \ \text{is determined by}\ 
  DF\vert^T_v\xi=D\pi\vert^T_v\eta
  \bigr\}. 
\end{align*}
%
%
Since $C_\mathcal{X}$ and $N^\ast(\partial{D_j})$ are Lagrangian submanifolds, we have 
%
%
\begin{align*}
  \operatorname{dim}\bigl(C_\mathcal{X}{\times}N^\ast({\partial}D_j)\bigr)
& = 
  \operatorname{dim}\bigl(\partial_-S(M)\bigr)
  +
  \operatorname{dim}(M^\text{int})
  +
  \operatorname{dim}(M^\text{int})
\\
& =
  (2n-2)+n+n
  =4n-2.
\end{align*} 
%
%
We set 
%
%
\begin{align}
  C
& :=
  \bigl(C_\mathcal{X}{\times}N^\ast({\partial}D_j)\bigr)\cap\mathcal{D}
\nonumber
\\
& =
  \bigl\{
  (\xi,\eta,\zeta) 
  \in 
  T^\ast\bigl(\partial_-S(M)\bigr){\times}T^\ast(M^\text{int}){\times}N^\ast(\partial{D_j}): 
\nonumber
\\
& \qquad 
  \exists v \in S(M^\text{int})\ \text{s.t.}\  
  \xi \in T^\ast_{F(v)}\bigl(\partial_-S(M)\bigr)\setminus\{0\}, 
\nonumber
\\
& \qquad
  \eta \in T^\ast_{\pi(v)}(M^\text{int})\setminus\{0\}, 
  \eta(v)=0, 
  DF\vert^T_v\xi=D\pi\vert^T_v\eta,\eta=\zeta\}
\nonumber
\\
& =
  \bigl\{
  (\xi,\eta,\eta): 
  \exists v \in S({\partial}D_j)\ \text{s.t.}\  
  \xi \in T^\ast_{F(v)}\bigl(\partial_-S(M)\bigr)\setminus\{0\},
\nonumber
\\
& \qquad 
  \eta \in N^\ast_{\pi(v)}(\partial{D_j})\setminus\{0\},  
  DF\vert^T_v\xi=D\pi\vert^T_v\eta\bigr\}.
\label{equation:intersection4-1}
\end{align}
%
%
Since $\xi$ is uniquely determined by $\eta$ and $v$ for $(\xi,\eta,\eta) \in C$, we deduce that 
%
%
$$
\operatorname{dim}(C)
=
\operatorname{dim}\bigl(S(\partial{D}_j)\bigr)
+
\operatorname{dim}\bigl(N^\ast_{\pi(v)}(\partial{D_j})\bigr)
=
(2n-3)+1
=
2n-2.
$$
%
%
The excess of $C$ is computed as  
%
%
\begin{align*}
  e
& =
  \operatorname{codim}(\mathcal{D})
  +
  \operatorname{codim}\bigl(C_\mathcal{X}{\times}N^\ast(\partial{D_j})\bigr)
  -
  \operatorname{codim}(C)
\\
& =
  \operatorname{codim}(\mathcal{D})
  -
  \operatorname{dim}\bigl(C_\mathcal{X}{\times}N^\ast(\partial{D_j})\bigr)
  +
  \operatorname{dim}(C)
\\
& =
  2n-(4n-2)+(2n-2)
  =
  0. 
\end{align*} 
%
%
If we show that $C$ is a clean intersection, then it follows that $C_\mathcal{X}{\circ}N^\ast(\partial{D_j})$  is transversal. 
Fix arbitrary $c=(\xi,\eta,\eta) \in C$, and pick up arbitrary 
$X=(X_1,X_2,X_3) \in T_c\bigl(C_\mathcal{X}{\times}N^\ast(\partial{D_j})\bigr)$. Consider a smooth curve on 
$C_\mathcal{X}{\times}N^\ast(\partial{D_j})$ of the form 
%
%
\begin{align*}
  \mathbb{R}\ni\alpha 
& \mapsto 
  \bigl(\xi(\alpha),\eta(\alpha),\zeta(\alpha)\bigr)
  =
  c+\alpha{X}+\mathcal{O}(\alpha^2) 
\\
& =
  (\xi+\alpha{X}_1,\eta+\alpha{X}_2,\eta+\alpha{X}_3)+\mathcal{O}(\alpha^2)
\end{align*}
%
%
near $\alpha=0$. Then we have 
%
%
$$
X
=
\frac{d}{d\alpha}\bigg\vert_{\alpha=0}\bigl(\xi(\alpha),\eta(\alpha),\zeta(\alpha)\bigr)
$$
%
%
If $X \in T_c(\mathcal{D})$, then we have 
%
%
$$
\frac{d}{d\alpha}\bigg\vert_{\alpha=0}\eta(\alpha)
=
\frac{d}{d\alpha}\bigg\vert_{\alpha=0}\zeta(\alpha),
$$
%
%
and 
%
%
$$
X=
\frac{d}{d\alpha}\bigg\vert_{\alpha=0}\bigl(\xi(\alpha),\eta(\alpha),\zeta(\alpha)\bigr)
=
\frac{d}{d\alpha}\bigg\vert_{\alpha=0}\bigl(\xi(\alpha),\eta(\alpha),\eta(\alpha)\bigr)
\in 
T_c(C). 
$$
%
%
Hence we obtain $T_c(\bigl(C_\mathcal{X}{\times}N^\ast(\partial{D_j})\bigr)){\cap}T_c(\mathcal{D}){\subset}T_c(C)$. 
Therefore we conclude that 
%
%
$$
T_c(\bigl(C_\mathcal{X}{\times}N^\ast(\partial{D_j})\bigr)){\cap}T_c(\mathcal{D})=T_c(C)
$$
%
% 
and that $C$ is clean intersection. 
Consider the projection 
$\pi_C:C\ni(\xi,\eta,\eta) \mapsto \xi \in T^\ast\bigl(\partial_-S(M)\bigr)$.  
It follows that $C$ is proper and $\pi_C^{-1}(\xi)$ for $\xi \in \pi_C(C)$ is connected since $\xi$ is uniquely determined by 
$\xi=(DF\vert^T_v)^{-1}{\circ}D\pi\vert^T_v$ for 
$\eta \in N^\ast_{\pi(v)}(\partial{D_j})$.    
%
%
\par
%
%
Next we show that $N^\ast(\Sigma_j)=C_\mathcal{X}{\circ}N^\ast(\partial{D_j})$. 
Using the expression \eqref{equation:intersection4-1} of $C$, we deduce that 
%
%
\begin{align}
  C_\mathcal{X}{\circ}N^\ast(\partial{D_j})
& =
  \bigl\{
  \xi: 
  \exists v \in S({\partial}D_j), \exists \eta \in N^\ast_{\pi(v)}(\partial{D_j})\setminus\{0\}, \text{s.t.}\
\nonumber
\\
& \qquad   
  \xi \in T^\ast_{F(v)}\bigl(\partial_-S(M)\bigr)\setminus\{0\},
  DF\vert^T_v\xi=D\pi\vert^T_v\eta\bigr\},
\label{equation:nstarsigmaj}
\\
  \Sigma_j
& :=
  \pi_{\partial_-S(M)}\bigl(C_\mathcal{X}{\circ}N^\ast(\partial{D_j})\bigr)
  =
  \{F(v) : v \in S(\partial{D_j})\}.
\nonumber 
\end{align}
%
%
Then we have $\operatorname{dim}(\Sigma_j)=\operatorname{dim}\bigl(S(\partial{D_j})\bigr)=2n-3$. 
Therefore we deduce that 
$\operatorname{codim}(\Sigma_j)=1$ in $\partial_-S(M)$ 
and that $\Sigma_j$ is a hypersurface in $\partial_-S(M)$. 
Since $C_\mathcal{X}$ is a homogeneous canonical relation and 
$N^\ast(\partial{D_j})$ is a conic Lagrangian submanifold, 
we conclude that $C_\mathcal{X}{\circ}N^\ast(\partial{D_j})$ is also a conic Lagrangian submanifold 
and that $N^\ast(\Sigma_j)=C_\mathcal{X}{\circ}N^\ast(\partial{D_j})$. 
This completes the proof. 
\end{proof}

\end{document}
